\documentclass[12pt]{article}

% --- Minimal, useful packages ---
\usepackage[margin=1in]{geometry}
\usepackage{amsmath,amssymb,amsthm,mathtools}
\usepackage{bm}            % bold math symbols
\usepackage{enumitem}      % compact lists
\usepackage{hyperref}      % clickable references
\usepackage{xcolor}        % (optional) for TODOs

% --- Theorem-like environments ---
\newtheorem{assumption}{Assumption}
\newtheorem{definition}{Definition}
\newtheorem{proposition}{Proposition}
\newtheorem{lemma}{Lemma}
\theoremstyle{remark}
\newtheorem{remark}{Remark}

% --- Handy macros ---
\newcommand{\E}{\mathbb{E}}
\newcommand{\Var}{\mathrm{Var}}
\newcommand{\Cov}{\mathrm{Cov}}
\newcommand{\Prob}{\mathbb{P}}
\newcommand{\1}{\mathbf{1}}
\newcommand{\argmax}{\mathop{\mathrm{arg\,max}}\nolimits}
\newcommand{\argmin}{\mathop{\mathrm{arg\,min}}\nolimits}
\newcommand{\dd}{\mathrm{d}}
\newcommand{\R}{\mathbb{R}}
\newcommand{\N}{\mathbb{N}}
\newcommand{\vect}[1]{\bm{#1}}
\newcommand{\set}[1]{\left\{\,#1\,\right\}}
\newcommand{\paren}[1]{\left(#1\right)}
\newcommand{\brak}[1]{\left[#1\right]}
\newcommand{\abs}[1]{\left|#1\right|}
\newcommand{\todo}[1]{\textcolor{red}{\textbf{TODO:} #1}}

% --- Title info ---
\title{Model Explanation - Mid/Later Life Interstate Moving}
\author{Tate Mason}
\date{\today}

\begin{document}
\maketitle

\begin{abstract}
  The purpose of this model is to estimate the motivations of different types of agents to move across states in the US. The model incorportates two sorts of agents,
  working and retired. The working agents are motivated by wage increases and the opportunity to buy a house. They give all time to labor, and at the end of each period have
  the opportunity to move. This is when they can purchase a house if their home savings are hight enough. The retired agents, however, are motivated by amenities 
  and non-monetary factors associated with moving. These agents do not work, of course, and still can move at the end of each period. Motivated by the data, these agents
  have a low probability of purchasing a home. Thus, I have taken away that choice and instead make it so in retirement, if a young agent did not buy a home, their home
  savings will meld into their retirement savings. The model, for simplicity, forgoes any sort of labor or health shocks, saying agents only have the trade-off of moving or not moving.
  The model is solved using value function iteration and the results will be compared to the data to see how well the model fits the data.
\end{abstract}

\section{Environment}
\begin{itemize}[leftmargin=1.2em]
  \item \textbf{Time:} Discrete $t=40, 41, ..., 85$. 
  \item \textbf{Agents:}
  \begin{enumerate}
    \item A continuum of working agents $\Pi_Y = \frac{1 -x}{(1-x) + (1-y)}$.
   \item A continuum of retired agents $\Pi_O = \frac{1 -y}{(1-x) + (1-y)}$.
  \item \textbf{Endowments/States:}
    \begin{itemize}
      \item Working agents: age $t$, home savings $a$, retirement savings $s$, income $w$, state $z$.
      \item Retired agents: age $t$, wealth $s$, pension $ss$, state $z$.
      \item \textbf{Uncertainty:} Probability of retirement each period $p^{retire}_t = (1-x)$, and probability of death $p^{die}_t = (1-y)$.
\end{itemize}

\section{Preferences}
    \subsection{Working Agents}
    Working agents have preferences over consumption $c$, housing $h(z)$, and state amentiies $m(z)$, and experience disutility from moving $\gamma(z')$:
    \begin{equation}
      U(c, h(z), m(z); t) = \frac{c^{1-\sigma}}{1 - \sigma} + \alpha_h \log(h(z)) + \alpha_m \log(m(z)) - \gamma(z') \cdot \1\{z' \neq z\}.,
    \end{equation}
    where $\sigma$ is the coefficient of relative risk aversion, $\alpha_h$ and $\alpha_m$ are the weights on housing and amenities,
    respectively, and $\gamma(z')$ is the moving cost to state $z'$.

    \subsection{Retired Agents}
    Retired agents, on the other hand only have preferences over consumption and amenities, and also experience disutility from moving:
    \begin{equation}
      U(c, m(z); t) = \frac{c^{1-\sigma}}{1 - \sigma} + \alpha_m \log(m(z)) - \gamma(z') \cdot \1\{z' \neq z\}.
    \end{equation}
    The parameters are the same as above, except there is no housing term. These agents do not have the option to buy a house, and thus do not derive utility from it.

\section{Information and Timing}
\begin{enumerate}[leftmargin=1.2em]
  \item At the beginning of each period, agents observe their state variables and the aggregate state $S_t$.
  \item Working agents choose consumption $c_t$ and investments $a_{t+1}$, $s_{t+1}$
  \item At the end of the period, agents will decide whether or not to move. If they move, they pay the moving cost $\gamma(z')$ and update their state to $z'$. 
  If they do not move, they remain in state $z$. They will then observe if their savings $a_t$ are high enough to purchase a house. If they can, they 
  will buy a house and their home savings will be reset to zero.
  \item Finally, they will observe if they retire.
  \item Retired agents choose consumption $c_t$ and investments $s_{t+1}$.
  \item At the end of the period, retired agents will decide whether or not to move. If they move, they pay the moving cost $\gamma(z')$ and update their state to $z'$. 
  If they do not move, they remain in state $z$. 
  \item Finally, they will observe if they die.
\end{enumerate}

If an agent dies, their wealth disappears from the economy. There is no bequest motive. They are then replaced by a newborn working agent with zero savings and randomly assigned state.

\section{Constraints and Market Clearing}
    \subsection{Working Agents}
      \begin{align}
        c_t + a_{t+1} + s_{t+1} \;=\; (1+r_t(z))a_t + (1+r^s_t(z))s_t + w_t(z) - \tau(w_t(z)) \\
        + (\1\{a_t \geq A_{house}\} \cdot (-A_{house}(z) + a_t))\1\{z' \ne z\},
      \end{align}
          where $A_{house}$ is the price of a house in state $z$, and $\tau(\cdot)$ is the fixed tax in state $z$. It is a 
          wasted tax, meaning it does not go to the government or anyone else in the economy. It is simply removed from the economy.

    \subsection{Retired Agents}
      \begin{equation}
        c_t + s_{t+1} \;=\; (1+r^s_t(z))s_t + ss_t(z) + \1\{a_t < A_{house}\} \cdot a_t,
      \end{equation}
          where the last term is the home savings that meld into retirement savings if the agent did not buy a house when they were working. In the next period after retiring, $s_{t+1} = s_t + a_t$.

Aggregate feasibility / market clearing:
\begin{equation}
  \int c_t(z) d\Pi_Y + \int c_t(z) d\Pi_O \;=\; \int w_t(z) d\Pi_Y + \int ss_t(z) d\Pi_O.
\end{equation}
The goods market clears, meaning all consumption must be equal to all income in the economy. There are no investments in this economy, so all income must be consumed.
There is no government spending, so the taxes are simply removed from the economy.

\section{Agent Problems}
\subsection{Working Agents}
\begin{equation}
  V_y(t, a, s, z) \;=\; \max\{V_y^m(t, a, s, z), V_y^s(t, a, s, z)\}, \\
\end{equation}
such that,
\begin{gather*}
  V^m_y(t, a, s, z) \;=\; \max_{c, a', s', z'} \; U(c, h(z), m(z); t) + x\beta \E[V_y(t+1, a', s', z')|z] \\
  \;+\; (1-x)\beta \E[V_o(t+1, 0, s', z')|z] \\
\end{gather*}
\begin{gather*}
  \quad \text{s.t.} \quad c + a' + s' = (1+r_t(z))a + (1+r^s_t(z))s + w_t(z) - \tau(w_t(z)) \\ 
  + (\1\{a \geq A_{house}\} \cdot (-A_{house}(z) + a))\1\{z'\ne z\}, \\
  \quad a' \geq 0, \; s' \geq 0, \; z' \in \{1, ..., Z\}, \\
\end{gather*}
\begin{gather*}
  V^s_y(t, a, s, z) \;=\; \max_{c, a', s'} \; U(c, h(z), m(z); t) + x\beta \E[V_y(t+1, a', s', z)|z] \\
  \;+\; (1-x)\beta \E[V_o(t+1, 0, s', z)|z] \\
\end{gather*}
\begin{gather*}
  \quad \text{s.t.} \quad c + a' + s' = (1+r_t(z))a + (1+r^s_t(z))s + w_t(z) - \tau(w_t(z)), \\
  \quad a' \geq 0, \; s' \geq 0.
\end{gather*}

State variables are time $t$, home savings $a$, retirement savings $s$, and state $z$. The value function is the maximum of moving and staying. If the agent moves, they choose consumption $c$,
next period home savings $a'$, next period retirement savings $s'$, and next period state $z'$. They receive utility from consumption, housing, and amenities, and pay a home purchase cost if they move. They then
face the probability of remaining working $x$ and receiving the working value function next period, or retiring with probability $1-x$ and receiving the retired value function next period. The agent faces the
budget constraint, which includes income, returns on savings, and the cost of moving and buying a house if they have enough home savings. If the agent does not move, they choose the same variables except for $z'$,
and do not have the option to buy a house.
 
\subsection{Retired Agents}
\begin{equation}
  V_o(t, s, z) = \max\{V_o^m(t, s, z), V_o^s(t, s, z)\}, \\
\end{equation}
such that,
\begin{gather*}
  V^m_o(t, s, z) = \max_{c, s', z'} \; U(c, m(z); t) + y\beta \E[V_o(t+1, s', z')|z] \\
\end{gather*}
\begin{gather*}
  \quad \text{s.t.} \quad c + s' = (1+r^s_t(z))s + ss_t(z) + \1\{a < A_{house}\} \cdot a,
  \quad s' \geq 0, \; z' \in \{1, ..., Z\}, \\
\end{gather*}
\begin{gather*}
  V^s_o(t, s, z) = \max_{c, s'} \; U(c, m(z); t) + y\beta \E[V_o(t+1, s', z)|z] \\
  \quad \text{s.t.} \quad c + s' = (1+r^s_t(z))s + ss_t(z) + \1\{a < A_{house}\} \cdot a,
  \quad s' \geq 0.
\end{gather*}

State variables are time $t$, retirement savings $s$, and state $z$. The value function is the maximum of moving and staying. If the agent moves, they choose consumption $c$,
next period retirement savings $s'$, and next period state $z'$. They receive utility from consumption and amenities, and pay a moving cost if they move. They then
face the probability of remaining retired $y$ and receiving the retired value function next period. The agent faces the budget constraint, which includes pension income, 
returns on savings, and the home savings that meld into retirement savings if they did not buy a house when they were working. If the agent does not move, they choose 
the same variables except for $z'$, and do not pay the moving cost.


\section{Conclusion}

This model provides a framework to understand the motivations of mid/later life agents to move across states in the US. By incorporating both working and retired agents, as well as
the option to purchase a house, the model captures the key trade-offs faced by these agents. The next step is to calibrate the model to match key moments in the data, such as the
proportion of movers by age, income, and presence of children. Once calibrated, the model can be used to conduct counterfactual experiments, such as changes in moving costs or housing prices,
to understand their impact on moving behavior. Overall, this model provides a useful tool to analyze mid/later life interstate moving.


\end{document}
